\subsection{Dynamic and static specification and invariant inference}
\label{subsec:rw3}

The pipeline proposed in this article depends on the ability to recover a draft formal contract from existing code, a capability with a long lineage in the specification- and invariant-inference literature. That literature divides naturally into dynamic and static traditions, and more recently into a learning-based strand; understanding where each is strong and where it is brittle clarifies precisely what our inference instrument contributes and, equally, what it deliberately does not claim.

The dynamic tradition observes a running program and reports properties that held over the executions seen. Its canonical instrument is Daikon, which instruments a target, records variable values at program points, and reports the templated properties (equalities, ordering, bounds, linear relationships) that survived every observed trace~\cite{ernst2007daikon}. Daikon established the central conceptual point that our brownfield characterisation phase reuses: a specification can be \emph{recovered} from behaviour rather than authored from scratch, lowering the bootstrap cost that has historically deterred adoption of formal contracts~\cite{meyer1992dbc}. The same conceptual move underlies specification mining, in which finite-state protocols governing API usage are learned from execution traces~\cite{ammons2002mining}. The defining limitation of every dynamic technique, however, is soundness: an inferred property is only \emph{likely} true, because it summarises the sampled executions, not all of them. This is exactly the distinction our architecture makes load-bearing. A code-derived clause is, in our terminology, a \texttt{code-inferred} obligation of \texttt{unchecked} assurance until OpenJML discharges it across all inputs, at which point it is promoted to \texttt{proved}. Dynamic inference supplies a hypothesis; it does not supply an oracle.

The static tradition derives invariants by sound over-approximation of program semantics. Abstract interpretation provides the foundational framework, computing fixpoints over abstract domains so that every reported property provably holds on every execution~\cite{cousot1977abstract}. Where abstract interpretation infers invariants from first principles, Houdini takes the complementary route of \emph{guess-and-check}: it postulates a large pool of candidate annotations and uses the ESC/Java verifier to refute the unprovable ones, retaining only the inductive remainder~\cite{flanagan2001houdini}. Houdini is the historical antecedent of the verifier-in-the-loop discipline our synthesis--verification loop generalises, and its candidate-refinement structure prefigures the counterexample-guided refinement (CEGIS) shape we adopt for code synthesis. The static methods buy soundness at the price of expressiveness and scale: rich abstract domains do not always converge, and the decidability and capacity limits of the underlying decision procedure bound what can be discharged. Our own honest reporting of Z3 timeouts on multiplicative reasoning, array theory, and rich preconditions is the contemporary form of precisely this limit, and it is the reason our tiered-assurance mechanism degrades full proof to bounded check to test rather than silently dropping an obligation.

A third, recent strand learns invariants and specifications directly. Data-driven frameworks such as ICE pose invariant synthesis as learning from examples, counterexamples, and implication pairs supplied by a verifier~\cite{garg2014ice}, and neural approaches reformulate the same problem as graph- or trace-conditioned prediction~\cite{si2018learning}. Since 2023 the focus has shifted to large language models (LLMs): models have been fine-tuned to predict Daikon-style invariants statically~\cite{pei2023llminvariants}, prompted to emit candidate loop invariants that symbolic tools then check~\cite{kamath2023loopy}, and augmented with learned re-ranking to triage many generated candidates against a sound verifier~\cite{chakraborty2023ranking}. These works share a structural commitment with our pipeline: the generative model is never trusted on its own, but is paired with a sound checker that admits only what it can prove. Yet they remain narrow---focused on the single artefact of a loop invariant for an isolated verification task---and they do not address provenance, the human-ratification problem, or evolution.

Across all three strands, the literature treats inference as a terminal deliverable: the inferred property is the output. Our contribution repositions inference as the \emph{opening move} of a verification-centric lifecycle. First, recovered clauses are not presented as fact but carried as provenance- and assurance-tagged obligations, so the near-circular nature of verifying code against a code-derived spec is made explicit rather than overclaimed as bug-freedom. Second, we triangulate the code-derived draft against two further intent streams---issue tickets and documentation---and treat \emph{disagreement} between streams as the signal of interest, surfacing candidate bugs that single-stream inference cannot. Third, because the inferred contracts are compositional, a change to a callee specification recomputes caller obligations and flags behavioural incompatibilities without execution, promoting our prior compositional-refinement result (Article~3 of this programme,~\cite{edmonds_article3}) from an inference technique to a lifecycle capability for version-compatibility analysis. No prior inference system combines recovery, multi-source reconciliation, sound verification across all inputs, and compositional change propagation into a single contract layer; that integration is the gap this article addresses.
